\subsection{A retained reference square for terminals $(1,2,c)$}
\label{proof:Q5-C184}
\begin{proposition}\label{prop:retained-reference-family}
Let $F_c=qf_2h_c$ in the canonical direction: the unit arm points
toward the terminal, the two-edge arm points outward, and the $c$-arm
has its first edge toward the terminal and its remaining edges outward.
The connector of length $L$ points from terminal to receiver.
For every outward receiving arm $a\ge3$, $c\ge4$, and $L\ge1$,
\begin{equation}\label{eq:Q184:scope}
 \E\bigl[E_of_a\{T^L(E_tF_c)\}^4\bigr]\le1-d/800.
\end{equation}
For the genuine receiving arm of length two, the same statement holds
when $c\ge5$:
\begin{equation}\label{eq:Q184:f2-extension}
 \E\bigl[E_of_2\{T^L(E_tF_c)\}^4\bigr]\le1-d/800.
\end{equation}
Here $E_t,E_o$ are the actual equal-pair factors at the two roots.
For $L\ge2$ both conclusions have the common stronger loss $d/600$.
\end{proposition}
\begin{proof}

\paragraph{The original measure and the two receiving masses.}
At $d=0$, $u=0$ and leaf integration gives equality. Assume $d>0$
for the proof.
Use $B=-K$, $u=B\one$, $q=1+u$, $p=1-u$ and the spectral measure
$\nu$ at $u$ from Section~\ref{sec:reference-common:receiving}.
It has mass $d$, support $[0,\theta]$, and $\int x^j\,d\nu=\delta_{j+1}$.
All means and adjoints below use the original finite measure.
For $a\ge3$ take
\[
 C_a=\one-G_a(-B)u-\beta_a\one\ge f_a,
 \quad\E C_a\le1-19d/16,\quad C_a\le p+6\sqrt d/5.
\]
For $a=2$ take exactly $C_2=f_2=p+Bu$, of mass $1-d$ and cap
$p+\sqrt d$. The proof below uses the weak coefficient-nine profile of Lemma~\ref{lem:reference-coefficient-nine}.

For every actual upper $Z=c_0\one+Q(B)u$, put
$\alpha=-\E Q(B)u$. The identities of
Lemma~\ref{lem:reference-common:receiving} read
\begin{equation}\label{eq:Q184:actual-moments}
 \operatorname{Var}Z=\int|Q|^2\,d\nu-\alpha^2,
 \qquad
 \operatorname{Cov}(C_a,Z)
 =-\int\operatorname{Re}(G_aQ)\,d\nu-\beta_{a+1}\alpha.
\end{equation}
They also hold at $a=2$, where $G_2=1+z$ and $\beta_3=d$.

\paragraph{One finite square and its positive continuation.}
For $c\ge5$, $d\le1/5$, use the actual $\sigma=9/4$ row of
Lemma~\ref{lem:reference-retained-packet}: $m=1-9d/8$, $k=m^{-1}$,
$U=(a-zb)u$, $R=-kU=(-t+zs)u$ and $D=2d/5$.
The retained coefficients keep $0\le t,s\le1$,
$t\ge k\tau+(1-x)s$, $a,b\ge x(1-x)$ and all signed $D_i$.
For $c=4$, $d\le1/4$, use instead
\[
 m=1-d,\quad U=(d+x-xz)u,\quad
 t=k(d+x),\quad s=kx,\quad D=5d/12.
\]
The actual arm sum is $f_2+h_4=2m-U$.
Its first-row error is bounded by $D$ in \eqref{eq:C169:6.1}.
Here $D_1=k(d^2+\delta_2)$, $D_2=k\delta_3$,
$D_3=-k(d\delta_2+\delta_3)$ and $D_4=-k\delta_4$.
These give $\gamma_3\ge0$, $\gamma_4\ge3d/4$ and
$\gamma_j\ge41d/96$ for $j\ge5$.
In both families
\begin{equation}\label{eq:Q184:duhamel}
 T^LF_c\le m^2Y_L,
 \quad
 Y_L=\one+G_{L+1}(B)u+B^LR+D\bar f_{L-1}
       -\sum_{j=3}^{L+1}\gamma_j\bar f_{L+1-j}.
\end{equation}
This is the finite identity of
Lemma~\ref{lem:reference-retained:continuation}, with the same direct
induction for the $c=4$ first row. Every discarded loss is nonnegative.

Retain no final loss at $L=1$, $\gamma_3$ at $L=2$,
both $g_3q$ and $g_4\one$ at $L=3$, where $g_i=\gamma_i$,
and $h=41d/96$ on each of the last two
paths at $L\ge4$. In every case the resulting upper is $W=m^2Z\ge T^LF_c$,
with the following full polynomial and unscaled mean:
\begin{equation}\label{eq:Q184:finite-uppers}
\begin{array}{c|l|l}
 L&Z=c_0\one+Q(B)u&\E Z\text{ upper}\\\hline
 1&c_0=1+D,\quad Q=1-xs+z(1-t)&1+D-d+D_1\\
 2&c_0=1+D-\gamma_3,\quad Q=Q_3+D&1-7d/20\\
 3&c_0=1+D-g_3-g_4,\quad Q=G_4-tz^3+sz^4+DG_2-g_3
    &1-31d/32\\
 \ge4&c_0=1+D-2h,\quad Q=G_{L+1}-tz^L+sz^{L+1}+DG_{L-1}-h
    &1-83d/96.
\end{array}
\end{equation}
The $L=2$ row is used here only for $c\ge5$; its exceptional
$c=4$ conclusion is \eqref{eq:Q167:T175.0}.
In that row $Q_3=1-x+tx+z(1-sx)$.
At $L=3$, $\alpha=\gamma_5+Dd\ge\tau+Dd$; at $L\ge4$,
$\alpha=\gamma_{L+2}+D\beta_L\ge h+D\tau$.
Thus the separate mean term in \eqref{eq:Q184:actual-moments}
has the favorable sign in all nonunit rows.
The genuine retained first loss has the uniform lower bound $g_3\ge d/8$:
for $c=4$, $D_1\le kd(d+\theta)\le7d/8$; for $c\ge5$,
the actual coordinate cap $a<16/25$ gives
$D_1<(40/31)(16/25)d=128d/155<7d/8$.
Also $0\le g_3\le d$ because $D_1\ge0$.
Thus $\E Z\le1+(5/12)d-d/8-2d/3-\kappa d=1-31d/32$.
At $L=3$ the retained polynomial subtracts $g_3$, while $g_4$ remains
a constant loss. This upper follows directly from
$\bar f_2=\one+G_2(B)u$ in the exact finite continuation.

\paragraph{Finite polynomial bounds.}
The finite alternating sums give real and normalized imaginary
coordinates in $[1-x,1]$ for $G_a$, $a\ge3$, and in
$[1-x,1-x+x^2]$ for $G_n$, $n\ge4$.
Consequently
\[
 \operatorname{Re}(G_aG_n)\ge0\ (n\ge4),\qquad
 \operatorname{Re}(G_aG_3)\ge1-3x+x^2\ge-1/16,\qquad
 \operatorname{Re}(G_aG_2)\ge1-2x\ge0.
\]
For $H_n=G_{n-1}+z^n$, the $n=4$ coordinates are
$(1-x+x^2,1)$; for $n\ge5$ they belong to the second rectangle above.
Thus $\operatorname{Re}(G_aH_n)\ge1-3x+2x^2-x^3>0$.
All these finite corners have real part at least $\kappa$.
They contract: $H_4,H_5$ have squared norms
$1-x+3x^2-2x^3+x^4$ and
$1-x-x^2+3x^3-2x^4+x^5$, respectively.
For $H_4$ the residual bracket increases to $-1451/32768$;
for $H_5$ use $x(3x-1)\le91/1024<1/4$ and
$-2x^4+x^5\le0$. For $n\ge6$ use
\[
 (1-z)H_n=1-z^{n-1}(1-z+z^2),\qquad
 |1-z+z^2|^2=1-x+x^2\le1,
 \quad 2\theta^{3/2}+\theta^4<1.
\]
The corresponding finite geometric identity gives $\|G_n\|\le1$
for $n=1$ and $n\ge3$; $\|G_2\|=\sqrt{1+\theta}<6/5$.
Hence the rectangle polynomial
$P_L=G_L-tz^{L-1}+sz^L$ contracts for every $L\ge4$.

The two shorter polynomials retain their sharper bounds
\begin{equation}\label{eq:Q184:tight-short-inputs}
 |Q_2|^2\le8/7,\qquad |Q_3|^2\le8/7\quad(c\ge5),
 \qquad Q_2=1-sx+z(1-t).
\end{equation}
These are \eqref{eq:retained-reference:Q2-tight} and
\eqref{eq:retained-reference:Q3-tight}, proved for the actual retained
coefficients and the full frequency interval. For $c=4$, $Q_3$ instead
has squared norm less than $7/6<(13/12)^2$.
Indeed its coordinates are affine in $k\in[1,4/3]$.
At $k=1$ the squared norm is at most one. At $k=4/3$ it is
\[
 1-x/3+(28/9)x^2-(40/9)x^3+(16/9)x^4+(16/9)x^5.
\]
Its second derivative, with $v=x-1/4$, is
$13/9-(28/3)v+48v^2+(320/9)v^3>0$:
for $v\le0$ use $v^3\ge-v^2/4$; for $v\ge0$ the quadratic
has minimum $107/108$. At $x=\theta$ its coordinates are
$243/256<19/20$, $599/768<4/5$, giving squared norm
less than $93/80<7/6$.

\paragraph{The weak-receiver rows and the genuine-$f_2$ correction.}
For $a\ge3$, \eqref{eq:Q184:actual-moments} gives the following
unscaled covariance budgets:
\begin{equation}\label{eq:Q184:covariance-budgets}
\begin{array}{c|c|c|c}
 L&\text{weak budget }b_0&\text{genuine }\operatorname{Cov}(f_2,Z)\text{ upper}
    &\E Z\text{ cap for transfer}\\\hline
 1&-3d/16&-439d/1024&9/8\\
 2&0&-3d/16&1\\
 3&d/16+d^2&-3d/16+d^2&1\\
 \ge4&29d^2/64&-3d/16+hd&1.
\end{array}
\end{equation}
For the unit weak row,
$\operatorname{Re}(G_aQ)\ge1-2x+x[t-(1-x)s]\ge1-2x$.
At $L=2$ the nonnegative coordinates and $a,b\ge x(1-x)$ give
$\operatorname{Re}(G_aQ)\ge1-3x+2x^2-x^3>0$.
At $L=3$ the $G_3$ corner costs $d/16$, and subtracting $g_3$ costs
at most $d^2$. At $L\ge4$ the subtracted $h$ costs $hd$ and
the possible $DG_3$ term costs at most $Dd/16$, giving $29d^2/64$.

For the genuine row use $G_2=1+z$ directly. At $L=1$,
\[
 \operatorname{Re}(G_2Q)=1-x+xt-sx
 \ge1-x-x^2\ge439/1024,\qquad \alpha=d-D_1\ge0.
\]
At $L=2$ its product is at least $1-2x\ge3/16$ and
$\alpha=d-D_2>0$.
At $L=3$ the corners are $G_3,G_4,G_5,H_4$; their genuine
products are at least $3/16$.
For longer connectors the $G$ corners give at least $\kappa^2$;
$H_5$ gives $1-2x+x^2-x^3>3/16$, and $H_n$, $n\ge6$, gives
at least $\kappa^2-\theta^3>3/16$.
The $D$ term and the separate $-d\alpha$ are favorable.
These facts prove the genuine column of the table.

Lemma~\ref{lem:reference-common:mass-transfer}, with $\eta=19/16$,
therefore replaces the genuine mass $1-d$ by the common comparison
mass $1-19d/16$ in the mean/covariance bracket. The required fees are
$27d/512$ at unit length and $3d/64$ otherwise; each is strictly
smaller than the gain in the table. In particular, the unscaled unit
mean may exceed one: it is at most $1+2d/5\le27/25<9/8$.
No change of measure or replacement of the actual receiving arm occurs.

\paragraph{Curvature on nonunit connectors.}
For $L\ge2$ the actual scaled upper $W=m^2Z$ satisfies
\[
 W\le(1-5d/8)q-d/2+(13/12)\sqrt d,\qquad
 \E W\le1-3d/2.
\]
Here the $L=2$ row uses $\|Q_2\|<15/14$ and
$\gamma_3\ge dz_0/m$, $z_0=19/32-89d/32>0$.
The actual $L=3$ pointwise upper is
\[
 W\le m^2[(1+D-g_3)q+(13/12+D)\sqrt d-g_4].
\]
Here $\|Q_3+D\|\le13/12+D$. To compare with the weak profile,
drop only the nonnegative loss $g_3q$, use $g_4\ge2d/3$ and
$D\le5d/12$. This gives $W\le m^2\mathcal B$, where
\[
 \mathcal B=(1+5d/12)q+(13/12+5d/12)\sqrt d-2d/3
       \ge(3/4)\sqrt d\ge0.
\]
The lower bound uses $q\ge0$ and $\sqrt d\le1/2$.
Scale this entire nonnegative row to get
$W\le(1-d)^2\mathcal B$. Its $q$ coefficient is at most
$(1-d)^2(1+5d/12)\le1-5d/8$.
Its remaining constant difference, with the real constant loss
$-2d/3$, is exactly $dR_3(\sqrt d)$ below.
For $L\ge4$ the row is bounded by
$(1-d)^2[(1-d/96)q+(1+d/2)\sqrt d-h]$.
After comparing the $q$ coefficients, the remaining constant differences
are $tR_2(t)$, $dR_3(t)$, $tR_4(t)$, respectively, where $t=\sqrt d$ and
\[
\begin{aligned}
 R_2(t)&=13/12-(15/14)(1-9t^2/8)^2
    +t[(1-9t^2/8)(19/32-89t^2/32)-1/2],\\
 R_3(t)&=1/6+(7/4)t-(4/3)t^2-t^3/4+(2/3)t^4-(5/12)t^5,\\
 R_4(t)&=1/12-(7/96)t+(3/2)t^2-(41/48)t^3
    +(41/96)t^5-t^6/2.
\end{aligned}
\]
Their Bernstein coefficients have minima $1/84,1/6,89/1152$
on $[0,9/20]$, $[0,1/2]$, $[0,1/2]$, respectively.
The scaled means follow from \eqref{eq:Q184:finite-uppers} and $m\le1-d$.
Thus the weak case of Lemma~\ref{lem:reference-coefficient-nine}
gives curvature nine for both receiving types.

The exact variance remains \eqref{eq:Q184:actual-moments} until bounded.
For $L=2$ it is at most
\[
 m^4d[8/7+2D+D^2]
 \le(8/7)d\exp[-15d/4-81d^2/32].
\]
This uses $\operatorname{Re}Q_3\le1$,
$\log m\le-9d/8-81d^2/128$, and
$1+(7/10)d+(7/50)d^2\le e^{(91/125)d}$.
For $L=3$, each full corner minus $g_3$ has norm at most $1-g_3/2$,
since its real part is at least $\kappa$ and
$2\kappa-1\ge3g_3/4$.
Adding $DG_2$ gives full norm at most
$1-g_3/2+(6/5)D\le1+7d/16$.
For $L\ge4$ the analogous two-loss norm is at most $1+13d/64$.
One may now discard the nonpositive variance mean-square to get
$d(1+7d/16)^2$ and $d(1+13d/64)^2$ before scaling by $m^4$.
For the former keep the quadratic logarithmic gain
$m^4\le(1-d)^4\le e^{-4d-2d^2}$, so its scaled variance is at most
$d e^{-25d/8-2d^2}$.

At $L=2$ the mean is bounded more sharply by
\[
 \E W\le m[1-(11027/5120)d+(9953/2560)d^2+d^4],
 \qquad \E W\le m^2(1-7d/20).
\]
Indeed $D_1+D_2\le k[\delta_2+\delta_3+(53/32)d^2
+(17/16)d\delta_2+d^4]$; its $E$ term is
$E[d^2-(1-d)\delta_2]\le d^4$.
For $d\le3/16$, comparison with the cubic lower Taylor polynomial
of $e^{-23d/8}$ leaves an upper difference at most
$d[-403/1000+(109/50)d-(2/5)d^2+d^3]<0$.
For $3/16\le d\le1/5$, the second upper has logarithm at most
$-(13/5)d-(53/40)d^2-(15/16)d^3\le-23d/8$.
Thus the three scaled mean bounds are
$e^{-23d/8}$, $e^{-95d/32}$ and $e^{-275d/96}$.
The common fourth expansion, with the signed bracket just established,
now gives
\[
 \begin{array}{c|c}
 L&\text{moment upper}\\\hline
 2&\mathcal U_3(d)\\
 3&e^{-209d/16}+9d e^{-25d/8-2d^2}
       +(d/4+4d^2)e^{-349d/32}=\mathcal U_4(d)\\
 \ge4&e^{-607d/48}+9de^{-115d/32}
       +(29/16)d^2e^{-339d/32}\le\mathcal U_5(d).
 \end{array}
\]
For the last inequality use $e>8/3$ and
$\sup_{s\ge0}s e^{-cs}=1/(ce)$.
Lemma~\ref{lem:retained-reference-payments} proves losses
$d/160$, $d/600$ and $d/60$ on the stated low domains.

\paragraph{Unit curvature and the two legal mean endpoints.}
For the actual unit row set $\mu=\E W$.
The relation $t\ge(1-x)s$ gives
$|(1-t)+zs|^2\le1$, hence
$W\le m^2(q+\sqrt d+D)$.
If $C_0\le p+6\sqrt d/5$, $m\le1-d$, $D\le5d/12$ and
$\mu\le m^2(1+2d/5)$, then
\begin{equation}\label{eq:Q184:unitnine}
 C_0[(W+\mu)^2+2\mu^2]\le9.
\end{equation}
To prove it, write $v=\sqrt d\le1/2$ and factor $m^4$.
The nonnegative polynomial upper is
$(p+6v/5)[(q+1+v+(49/60)v^2)^2+2(1+(2/5)v^2)^2]$.
Its coefficients are bounded by
$13/2,66/5,65/4,25/2,4,6/5$.
The terms of degrees at least two are at most $24v^2$.
Comparing with $9e^{4v^2}$ reduces to
$P(v)=5/2-(66/5)v+12v^2+72v^4>0$.
On $[0,1/4]$ it decreases to $37/160$; afterwards
$P''\ge78$, $P'(1/4)=-27/10$, giving
$P\ge37/160-243/5200>0$.
The actual unit means obey the required cap: for $c=4$,
$D_1\le7d/8$ gives $\mu/m^2\le1+7d/24$;
for $c\ge5$, $D_1<d$ gives $\mu/m^2\le1+2d/5$.

In the mean/covariance bracket after the genuine-mass compensation,
use $h_0=1-19d/16$ and
$f(y)=h_0y^4-(3/4)m^2dy^3$.
Its maximum on $[0,M]$ is zero or $f(M)$.
The following actual scaled mean bounds have $M>1/2$ and $f(M)>0$:
\[
\begin{array}{c|c|cc}
 c&M&e^{-a d}\le M&M\le e^{-\lambda d}\\\hline
 4&(1-d)[1-(113/96)d+(19/12)d^2]&a=209/96&\lambda=129/64\\
 \ge5&(1-9d/8)[1-(211/160)d+(373/160)d^2]&a=391/160&\lambda=35/16.
\end{array}
\]
These follow by bounding $D_1$ by $k(d^2+\delta_2)$ and
$k[\delta_2+(53/32)d^2]$, respectively.
Write each cubic as $M=1-ad+bd^2-hd^3$.
The log-convexity numerator is
\[
 MM''-(M')^2=(2b-a^2)+(2ab-6h)d-2b^2d^2+4bhd^3-3h^2d^4.
\]
At $(a,b,h,D)=(209/96,265/96,19/12,1/4)$ its Bernstein
coefficients are
\[
 7199/9216,\quad7689/8192,\quad414635/442368,
 \quad41561/49152,\quad5753/8192;
\]
at $(391/160,4883/1280,3357/1280,1/5)$ they are
\[
\begin{gathered}
 42439/25600,\quad3693013/2048000,\quad215610671/122880000,\\
 163026843/102400000,\quad1395197823/1024000000.
\end{gathered}
\]
They are positive. The tangent at zero gives the lower exponential;
the endpoint chord gives the upper, using
\[
 M(1/4)=309/512<(2329/3840)(255/256)<e^{-129/256},
\]
\[
 M(1/5)=51429/80000<1-7/16+(7/16)^2/2-(7/16)^3/6<e^{-7/16}.
\]
The first comparison uses the degree-five lower Taylor polynomial
for $e^{-1/2}$ and $e^{-1/256}\ge255/256$.
Since $h_0\ge45/64$, $M>1/2$ and $(3/4)m^2d\le3/16$,
the positive endpoint is valid before replacing its mean factors.

For $c=4$ the nonnegative unit coordinates decrease with
$k\ge1$, so
$|Q|^2\le(1-x^2)^2+x(1-x)^2\le1+x-3x^2\le13/12$.
Moreover
$\alpha=d-D_1\ge d(\kappa-2d)/m>0$. Retaining the negative square gives
\[
 \operatorname{Var}W\le(13/12)m^4d-m^2(\kappa-2d)^2d^2.
\]
For $c\ge5$, $a,b\ge x(1-x)$ give
\[
 |Q|^2\le1+x-4x^2+5x^3-x^4-x^5+x^6
       \le1+x-(13/6)x^2.
\]
The last inequality follows since
$4-5x+x^2+x^3-x^4\ge2279183/1048576>13/6$.
Using $\delta_3\ge\delta_2^2/d$ and maximizing the resulting quadratic
gives $\int|Q|^2\,d\nu\le29d/26$.
The inequalities $m^2\ge e^{-7d/3}$ and $m^2\ge e^{-8d/3}$,
respectively, together with the two valid mean bounds, yield moment
uppers $\mathcal U_1(d)$ and $\mathcal U_2(d)$.
In the second row the negative cubic term is at least
$(3/4)d e^{-10d}$, because $8/3+3(391/160)=4799/480<10$.
Lemma~\ref{lem:retained-reference-payments} proves losses $d/800$
and $d/16$ for the unit low domains.

\paragraph{A single high comparison in the same direction.}
For $c\ge5$, \eqref{eq:compact-strong-path-means} gives, throughout
$d\ge1/5$,
\[
 h_c\le2\mu_{c-1}\le2-(83/16)d\le77/80,
 \qquad T^LF_c\le(77/80)T^L(qf_2).
\]
The genuine-$f_2$ reference for $qf_2$ is proved in
Section~\ref{proof:Q5-C032} at $L=1$,
Section~\ref{proof:Q5-C037} at $L=2,3,4,5$, and
Section~\ref{proof:Q5-C033} at $L\ge6$.
For $a\ge3$, the actual all-length proof comprises
Section~\ref{proof:Q5-C187} at $L=1$,
\eqref{eq:Q191:3.1} at $L=2$,
\eqref{eq:Q191:11.1} at $L=3,4,5$ and
\eqref{eq:q117-unified:shortone-two-tail} at $L\ge6$.
Each is at most one. Thus all these fourth moments are at most
$(77/80)^4<1-1/2400\le1-d/800$; the cap also covers the stated
stronger nonunit loss.

For $c=4$, $d\ge1/4$, one has $h_4\le2(1-2d)\le1$.
At $L=2$ use \eqref{eq:Q167:T175.0}, with loss $d/16$.
At $L=3,4,5$ the same-direction $qf_2$ bound
\eqref{eq:Q191:11.1} has loss $9d/10$; for $L\ge6$,
\eqref{eq:q117-unified:shortone-two-tail} has loss $3d/800$.
All exceed the required $d/600$.
At unit length, the actual unpropagated square gives
\[
 \E F_4\le H(d)=1-2d+[-2+\theta(1+\theta)/2]d^2+(9/2)d^3.
\]
This follows from
$\E F_4\le m^2-m(d^2+2\delta_2)
+(d^3+2d\delta_2+\delta_3+\delta_4)/2$.
The moment bounds imply the displayed $H$. Its derivative increases
on $[1/4,1/3]$ and has $H'(1/3)<-3/2$, while $H(1/4)<15/32$.
Consequently $TF_4<15/16$ and its actual reference fourth is less
than $15/16<1-d/800$.

Finally positivity inserts $0\le E_t,E_o\le1$ into the proved
inequalities.
All arm and connector directions are fixed throughout the host domain.
\end{proof}
